%% lahu-master.tex
%% Hand-authored scaffold for The Dictionary of Lahu, TEI Lex-0 digital
%% edition. Owns everything the generated dictionary body needs but
%% doesn't itself contain: documentclass/packages, page geometry, the
%% body font, the entry-formatting macros (\headword, \dividerletter,
%% etc.), the running footer, the title page, and the table of
%% contents. The dictionary body itself -- generated/latex/lahu.tex,
%% produced from generated/tei/lahu.xml by src/lahu-latex.xsl -- is
%% \input below, inside a multicols block this file opens and closes.
%%
%% Compile this file (not generated/latex/lahu.tex) with XeLaTeX,
%% always from the project root so the \input paths below resolve;
%% see README.md / README-tei.md for the exact command, or just run
%% src/regenerate-all-files.sh which does this for you:
%%
%%   xelatex -interaction nonstopmode -output-directory=generated/latex \
%%     -jobname=lahu src/latex/lahu-master.tex
%%   xelatex -interaction nonstopmode -output-directory=generated/latex \
%%     -jobname=lahu src/latex/lahu-master.tex     (again, for TOC/page refs)
%%
%% -jobname=lahu keeps the output named generated/latex/lahu.pdf (etc.)
%% rather than lahu-master.pdf, matching the pre-scaffold layout.
%%
%% Front/back matter: this file is the place to add it. \input a new
%% .tex file here (before the dictionary body for front matter, after
%% it for back matter) the same way title.tex and toc.tex are below.

\documentclass[10pt,twoside]{article}

%% Packages
\usepackage{fontspec}
\usepackage{multicol}
\usepackage{geometry}
\usepackage{microtype}
\usepackage{xcolor}
\usepackage{hyperref}
\usepackage{enumitem}
\usepackage{fancyhdr}
\usepackage{newunicodechar}

%% Page geometry
\geometry{
  a4paper,
  top=2cm, bottom=2.5cm,
  inner=2cm, outer=1.5cm,
  headsep=0.5cm
}

%% Body font. Chosen for broad Unicode coverage of IPA Extensions and
%% combining diacritical marks (Lahu's tone marks are combining accents
%% over vowels and consonants, e.g. g with combining diaeresis for g̈),
%% and because it ships with most complete TeX Live installations. If
%% XeLaTeX reports this font isn't found, either install DejaVu Serif
%% or change the name below to any Unicode font you have that covers
%% combining diacritics well, e.g. "Charis SIL" (built specifically for
%% this kind of linguistic data; see https://software.sil.org/charis/)
%% or "Doulos SIL".
%% NB: deliberately NOT Ligatures=TeX. That option treats a bare `'`/`"`
%% as a CLOSING curly quote (TeX's classic convention expects an opening
%% backtick/`` `` for the other side); since our source only ever emits
%% plain straight `'`/`"` on both sides of a quotation (never backticks),
%% every one of them was coming out as a right-facing curly quote at
%% render time (’.../’..., not ‘.../’...) even though the underlying text
%% is plain ASCII throughout the whole pipeline. Confirmed by extracting
%% the compiled PDF's text layer. We don't need TeX's dash-ligation
%% either -- em dashes already arrive as real U+2014 characters straight
%% out of dl_convert.py, never as literal "--". See also latex-escape in
%% lahu-latex.xsl, which now emits \textquotesingle{}/\textquotedbl{} for
%% '/" as a second, independent guard against this regardless of font
%% ligature settings.
\setmainfont{DejaVu Serif}[Renderer=OpenType]

%% Monospace font, used only by the front/back-matter documents
%% generated by src/process_front_and_back_matter.py (see
%% generated/latex/frontmatter.tex, backmatter.tex): several of those
%% originals (lists, glossaries, the abbreviation tables) used a fixed-
%% width font with space-aligned columns, which \ttfamily reproduces
%% via that script's monospace rendering mode. Without an explicit
%% \setmonofont, fontspec's \ttfamily falls back to Latin Modern Mono,
%% which lacks the IPA Extensions / Lahu diacritics used throughout
%% this corpus -- confirmed by a wall of "Missing character" warnings
%% for ɛ/ɔ/ʔ/ɨ/etc. before this was added. DejaVu Sans Mono has full
%% coverage of everything those documents need except the two
%% standalone modifier-letter accents (next block) and ⪤ (already
%% handled below), and ships alongside DejaVu Serif in the same
%% font family/install.
\setmonofont{DejaVu Sans Mono}[Renderer=OpenType]

%% Per-character fallback fonts for symbols DejaVu Serif doesn't cover.
%% "DejaVu Math TeX Gyre" ships alongside DejaVu Serif in any reasonably
%% complete TeX Live install (it's part of the same dejavu/tex-gyre-math
%% font collection), so this should be available wherever DejaVu Serif
%% itself is; \IfFontExistsTF guards it anyway so a missing font here
%% degrades to a normal "Missing character" warning instead of a hard
%% compile error.
\IfFontExistsTF{DejaVu Math TeX Gyre}{%
  \newfontfamily\dejavumathfallback{DejaVu Math TeX Gyre}[Renderer=OpenType]
  \newunicodechar{≣}{{\dejavumathfallback ≣}}% U+2263 STRICTLY EQUIVALENT TO
}{}

%% ⪤ (U+2AA4 GREATER-THAN OVERLAPPING LESS-THAN, our "cognate with"
%% connector -- see the '&' entry in LAHU_UNICODE in dl_convert.py) is
%% NOT in DejaVu Serif, DejaVu Math TeX Gyre, or any other font in a
%% standard TeX Live install -- confirmed by scanning every font's
%% cmap. "STIX Two Math" does have it. This is not part of a base TeX
%% Live install, so it needs a one-time manual install; see README.md
%% for the install command. As with the DejaVu Math TeX Gyre fallback
%% above, \IfFontExistsTF guards it: if the font isn't installed, ⪤
%% degrades to a normal "Missing character" warning at compile time
%% (a visible blank/box glyph in the PDF) rather than a hard error.
\IfFontExistsTF{STIX Two Math}{%
  \newfontfamily\stixmathfallback{STIX Two Math}[Renderer=OpenType]
  \newunicodechar{⪤}{{\stixmathfallback ⪤}}% U+2AA4 GREATER-THAN OVERLAPPING LESS-THAN
}{}

%% ˊ/ˋ (U+02CA/U+02CB, MODIFIER LETTER ACUTE/GRAVE ACCENT -- the
%% "bare"/standalone form of a citation-form diacritic that isn't
%% composed onto a base letter; see STANDALONE_DIAC in dl_convert.py)
%% are in DejaVu Serif but NOT in DejaVu Sans Mono -- surfaced only
%% once the front/back-matter monospace tables above started using
%% these characters (e.g. BIBLIOG.TXE's bare accents on a journal
%% abbreviation). No \IfFontExistsTF guard needed here: DejaVu Serif
%% is the required \setmainfont itself, not an optional fallback font.
\newunicodechar{ˊ}{{\rmfamily ˊ}}
\newunicodechar{ˋ}{{\rmfamily ˋ}}

%% Dictionary entry macros
%%
%% Hanging indent, two levels deep, via \leftskip/\hangindent rather
%% than a fixed \hspace on the first line: an entry's own wrapped lines
%% indent by \dictindent (the column where sub-entries' headwords
%% start); a sub-entry's first line starts at \dictindent, and ITS
%% wrapped lines indent one further \dictindent, to 2*\dictindent. Each
%% macro below sets \leftskip and \hangindent itself at the start of
%% its own paragraph, so nothing leaks from one paragraph to the next.
%%
%% \headword{Lahu}              -- bold Lahu headword (top level)
%% \graminfo{text}              -- grammatical information in small italic
%% \usgetym{text}               -- loanword flag
%% \usglabel{text}              -- bracketed register/dialect/cross-ref label
%% \sensenum{n}                 -- sense number
%% \examplecit{Lahu}{Eng}       -- example + translation (top level; translation
%%                                  is not quoted, just set off by a quad space;
%%                                  Lahu text set 2pt smaller than body text;
%%                                  indented one level deeper than the entry
%%                                  it appears with, starts at \dictindent,
%%                                  wraps to 2*\dictindent, same as \notetext)
%% \notetext{text}              -- lexicographic note (top level), rendered
%%                                  "/ text /", indented one level deeper than
%%                                  the entry it annotates (starts at
%%                                  \dictindent, wraps to 2*\dictindent)
%% \subentry{Lahu}{content}     -- sub-entry headword (bold, 2pt smaller than
%%                                  \headword) + gram + senses
%% \subexamplecit{Lahu}{Eng}    -- example + translation (inside a sub-entry;
%%                                  Lahu text set 2pt smaller, as \examplecit;
%%                                  indented one level deeper than the
%%                                  sub-entry, starts at 2*\dictindent, wraps
%%                                  to 3*\dictindent, same as \subnotetext)
%% \subnotetext{text}           -- lexicographic note (inside a sub-entry),
%%                                  "/ text /", indented one level deeper than
%%                                  the sub-entry it annotates (starts at
%%                                  2*\dictindent, wraps to 3*\dictindent)
%% \editorialnote{text}         -- pipeline diagnostic (gray, review mode only)
%% \dividerletter{x}            -- letter-divider (from the original "PUT x
%%                                  HERE" typesetting instructions): closes
%%                                  the current multicols block, drops a
%%                                  \label{div:x} for the table of contents
%%                                  (see toc.tex) to \pageref, sets the
%%                                  running head (\markright, shown via
%%                                  \rightmark on odd/recto pages -- see
%%                                  the header setup below) to this letter,
%%                                  prints a full-width ornamental
%%                                  rule/letter/rule, then reopens multicols
%%                                  with the same
%%                                  \raggedright/\small/\dictindent state as
%%                                  the top of the document, since closing
%%                                  multicols ends the group those were set
%%                                  in. Safe to invoke from anywhere in the
%%                                  entry stream (even mid sub-entry) because
%%                                  the XSLT template that calls it emits it
%%                                  as a flat run of tokens, never inside an
%%                                  open brace group.
%%
%%                                  IMPORTANT: this uses plain \label, not
%%                                  \addcontentsline. \addcontentsline here
%%                                  (tried first) corrupted one page of the
%%                                  compiled PDF into blank/wrong content
%%                                  when all 33 dividers were present -- a
%%                                  real interaction bug between hyperref's
%%                                  \addcontentsline and multicol's deferred
%%                                  page-breaking, reproduced and confirmed
%%                                  by removing it; see README-tei.md. Do
%%                                  not reintroduce \addcontentsline /
%%                                  \tableofcontents for this without
%%                                  re-testing a full 662-page compile.
%%
%% Declared here (must be in the preamble); the actual width is set
%% in \small below, once we're in the font size it will be used in --
%% "em" is resolved at the point \setlength runs, not dynamically.
\newlength{\dictindent}

\newcommand{\headword}[1]{%
  \par\vspace{2pt}%
  \leftskip=0pt \hangindent=\dictindent \hangafter=1
  \noindent\textbf{#1}\enspace%
}
\newcommand{\graminfo}[1]{%
  {\small\textit{#1}}\enspace%
}
\newcommand{\usgetym}[1]{%
  {\small\textit{[#1]}}\enspace%
}
\newcommand{\usglabel}[1]{%
  {\small\textit{[#1]}}\enspace%
}
\newcommand{\sensenum}[1]{%
  \textbf{#1.}\enspace%
}
\newcommand{\examplecit}[2]{%
  \par
  \leftskip=\dictindent \hangindent=\dictindent \hangafter=1
  \noindent{\fontsize{7pt}{8.4pt}\selectfont\textit{#1}}\quad #2%
}
\newcommand{\notetext}[1]{%
  \par
  \leftskip=\dictindent \hangindent=\dictindent \hangafter=1
  \noindent / #1 /%
}
\newcommand{\subentry}[2]{%
  \par\vspace{1pt}%
  \leftskip=\dictindent \hangindent=\dictindent \hangafter=1
  \noindent{\fontsize{7pt}{8.4pt}\selectfont\textbf{#1}}\enspace#2%
}
\newcommand{\subexamplecit}[2]{%
  \par
  \leftskip=2\dictindent \hangindent=\dictindent \hangafter=1
  \noindent{\fontsize{7pt}{8.4pt}\selectfont\textit{#1}}\quad #2%
}
\newcommand{\subnotetext}[1]{%
  \par
  \leftskip=2\dictindent \hangindent=\dictindent \hangafter=1
  \noindent / #1 /%
}
\newcommand{\editorialnote}[1]{%
  \par{\small\color{gray}[Editorial note: #1]}%
}
\newcommand{\dividerletter}[1]{%
  \end{multicols}%
  \label{div:#1}%
  \markright{#1}%
  \begin{center}
    \rule{0.35\linewidth}{0.4pt}\quad
    {\fontsize{28pt}{28pt}\selectfont\bfseries\itshape #1}%
    \quad\rule{0.35\linewidth}{0.4pt}
  \end{center}
  \begin{multicols}{2}
  \raggedright
  \small
  \setlength{\dictindent}{1.3em}%
}

%% Column separator
\setlength{\columnsep}{1.2em}
\setlength{\columnseprule}{0.4pt}

%% Hyperref
\hypersetup{colorlinks=true, linkcolor=black, urlcolor=blue}

%% Running head/footer for body pages. \thispagestyle{empty} (used on
%% the title page and TOC page) suppresses both there.
%%
%% Footer: page number in the outer corner (left on even/verso pages,
%% right on odd/recto pages, standard twoside book convention), the
%% Lahu collation sequence (alphabetization order, reconstructed from
%% the HEADER file in the original WordStar source; see README.md)
%% centered.
%%
%% Header: a stub running head, alternating by page side --
%% the book title on even/verso (left) pages, the current letter
%% section on odd/recto (right) pages. The current-section text comes
%% from \rightmark, which \dividerletter (below) updates via
%% \markright each time a new letter section starts. NOTE: \markright
%% inside a multicols block is a known rough edge (multicol's deferred,
%% per-page output routine doesn't always see mark changes exactly
%% when you'd expect); this was tested with a full 33-divider, 662-page
%% compile and came out correct, but if a future edit ever makes
%% \rightmark look stale or one-section-behind on some pages, that's
%% the mechanism to suspect first -- see the \dividerletter comment
%% above and the \addcontentsline bug note there for the related,
%% worse-behaved thing this project already hit once.
\pagestyle{fancy}
\fancyhf{}
\renewcommand{\headrulewidth}{0.2pt}
\renewcommand{\footrulewidth}{0.2pt}
\fancyhead[LE]{\slshape The Dictionary of Lahu}
\fancyhead[RO]{\slshape\rightmark}
\fancyfoot[LE,RO]{\thepage}
\fancyfoot[C]{\scriptsize a\ á\ â\ à\ ā\ âʔ\ àʔ\ \textbar\ a\ i\ u\ e\ o\ ɛ\ ɔ\ ɨ\ ə\ \textbar\ q\ qh\ k\ kh\ g\ ŋ\ c\ ch\ j\ t\ th\ d\ n\ p\ ph\ b\ m\ h\ g̈\ š\ y\ f\ v\ l}

\begin{document}

%% Front matter (title page, recovered dedication/preface/plates-list/
%% acknowledgments/symbols-and-abbreviations, and the hand-built letter
%% table of contents) is roman-numbered, matching standard book
%% convention and the original's own scheme (CONTENTS.TXE numbers the
%% Press-supplied preliminaries i-viii) -- reset to arabic 1 at the
%% dictionary body below. generated/latex/frontmatter-pre-toc.tex and
%% frontmatter-post-toc.tex are produced by
%% src/process_front_and_back_matter.py, not hand-authored; see that
%% script's module docstring for which originals/DLOtherFiles/ files
%% are included, in what order, and why (split into "pre-toc"/
%% "post-toc" because the original book's own front matter interleaves
%% the dedication before its table of contents and the rest --
%% acknowledgments, symbols and abbreviations, etc. -- after it, per
%% CONTENTS.TXE). Each document \input from these two files begins
%% with its own \clearpage, so nothing further is needed between them.
\pagenumbering{roman}

%% \pagestyle{fancy} (declared above) is for the dictionary body only --
%% its footer is a navigational aid specific to the alphabetical
%% dictionary (the collation-sequence strip) and doesn't belong on
%% front/back matter. Switch to the plain style (centered page number,
%% no header/footer strip) for the whole front-matter section here;
%% each generated document's own \thispagestyle{plain} (see
%% process_front_and_back_matter.py's render_document) only covers that
%% document's OWN first page, so a persistent \pagestyle change is
%% still needed for any document that runs past one page (confirmed
%% necessary: Acknowledgments' second page fell back to the fancy
%% dictionary footer without this).
\pagestyle{plain}

%% title.tex -- standalone title page for The Dictionary of Lahu.
%% \input from src/latex/lahu-master.tex, before the table of contents
%% and the dictionary body. Edit this file directly to change the
%% title page; it is hand-authored, not generated.

\begin{titlepage}
\centering
\vspace*{4cm}

{\Huge\bfseries The Dictionary of Lahu}\par
\vspace{1.2cm}
{\Large James A. Matisoff}\par
\vspace{0.2cm}
{\large 1988}\par

\vspace{2.5cm}
{\large TEI Lex-0 Digital Edition}\par

\vfill
{\small Transduced from a Lexware (Shoebox/MDF-style) database
  reconstructed from the original WordStar/CP-M source files.\par}

\vspace{1cm}

\end{titlepage}

\clearpage

\input{generated/latex/frontmatter-pre-toc.tex}
\clearpage
%% toc.tex -- table of contents for The Dictionary of Lahu.
%% \input from src/latex/lahu-master.tex, after the front matter that
%% precedes it (title page, dedication, preface in Lahu) and before the
%% front matter that follows it (list of plates, acknowledgments,
%% symbols and abbreviations, introduction, dictionary body, back
%% matter) -- see lahu-master.tex's own \input sequence, and
%% src/front-back-matter-order.csv for the reviewed, approved order
%% this file's entries follow. Matches the order and section numbering
%% of originals/DLOtherFiles/CONTENTS.TXE's own "TABLE OF CONTENTS"
%% listing (lines 49-123), plus two back-matter appendices CONTENTS.TXE
%% doesn't mention (Lahu Linguistic Terminology, Vocabulary Supplement)
%% that are nonetheless included -- see process_front_and_back_matter.py's
%% module docstring for why.
%%
%% This is a hand-built list, not a \tableofcontents.
%% src/process_front_and_back_matter.py drops a \label{toc:LABEL} at
%% the top of each front/back-matter document it generates (LABEL
%% matches the 'label' field in that script's manifest -- the two files
%% must be kept in sync by hand if either changes); \dividerletter (see
%% lahu-master.tex) drops \label{div:X} at each of the 33 letter-section
%% breaks in the dictionary body. \tocline below looks either kind of
%% label up with \pageref, so front-matter entries automatically show
%% roman numerals and dictionary-body/back-matter entries automatically
%% show arabic numerals -- \pageref just echoes whatever \pagenumbering
%% was active on the page a label sits on; no separate handling needed
%% here for that distinction.
%%
%% \tocsub is for outline sub-entries with no page of their own to link
%% to (the Introduction's missing 1.0/2.0/2.1/3.0/4.0-4.5, and the
%% Bibliography's I/II/III, which are internal headings within one
%% continuous document rather than separately paginated sections).
%%
%% \tableofcontents + \addcontentsline was tried first and rejected:
%% with all 33 dividers present, \addcontentsline corrupted one page
%% of the compiled PDF (a genuine hyperref/multicol interaction bug
%% with anchors dropped right at a multicols close/reopen boundary
%% inside a large multicols block) -- reproduced, confirmed by
%% removing it, and documented in lahu-master.tex's \dividerletter
%% comment and README-tei.md. Plain \label/\pageref doesn't hit it.
%%
%% Needs a second XeLaTeX pass (see lahu-master.tex's compile command)
%% for page numbers to be correct, same as any LaTeX cross-reference.

\newcommand{\tocline}[2]{\noindent #1\dotfill\pageref{#2}\par}
\newcommand{\tocsub}[1]{\noindent\hspace{1.5em}#1\par}

\thispagestyle{empty}
\begin{center}
{\Large\bfseries Contents}
\end{center}
\vspace{1.5em}

\tocline{List of Plates}{toc:list-of-plates}
\tocline{Acknowledgments}{toc:acknowledgments}
\tocline{Symbols and Abbreviations}{toc:symbols-and-abbreviations}
\tocline{Introduction to the \textit{Dictionary}}{toc:introduction}
\tocsub{1.0\quad History of the Lahu dictionary project}
\tocsub{2.0\quad The Lahu people and the Lahu language}
\tocsub{\hspace{1.5em}2.1\quad The genetic position of Lahu}
\tocline{\hspace{1.5em}2.2\quad Lahu dialects and cultural subdivisions}{toc:intro-2-2}
\tocsub{3.0\quad Transcription and alphabetical order of entries}
\tocsub{4.0\quad Structure and format of the individual entry}
\tocsub{\hspace{1.5em}4.1\quad Lemmata and lemmatizational dilemmas}
\tocsub{\hspace{1.5em}4.2\quad Form-classes}
\tocsub{\hspace{1.5em}4.3\quad Glosses}
\tocsub{\hspace{1.5em}4.4\quad Subentries}
\tocsub{\hspace{1.5em}4.5\quad Remarks: cross references and etymologies}

\vspace{0.8em}
\tocline{The Dictionary of Lahu}{div:a}

\begin{multicols}{2}
\raggedright
\tocline{a}{div:a}
\tocline{i}{div:i}
\tocline{u}{div:u}
\tocline{e}{div:e}
\tocline{o}{div:o}
\tocline{ɛ}{div:ɛ}
\tocline{ɔ}{div:ɔ}
\tocline{ɨ}{div:ɨ}
\tocline{ə}{div:ə}
\tocline{q}{div:q}
\tocline{qh}{div:qh}
\tocline{k}{div:k}
\tocline{kh}{div:kh}
\tocline{g}{div:g}
\tocline{ŋ}{div:ŋ}
\tocline{c}{div:c}
\tocline{ch}{div:ch}
\tocline{j}{div:j}
\tocline{t}{div:t}
\tocline{th}{div:th}
\tocline{d}{div:d}
\tocline{n}{div:n}
\tocline{p}{div:p}
\tocline{ph}{div:ph}
\tocline{b}{div:b}
\tocline{m}{div:m}
\tocline{h}{div:h}
\tocline{g̈}{div:g̈}
\tocline{š}{div:š}
\tocline{y}{div:y}
\tocline{f}{div:f}
\tocline{v}{div:v}
\tocline{l}{div:l}
\end{multicols}

\vspace{0.8em}
\tocline{Appendix: Lahu Linguistic Terminology}{toc:appendix-lingterm}
\tocline{Appendix: Vocabulary Supplement --- Birdnames}{toc:appendix-birdlist}
\tocline{Bibliography}{toc:bibliography}
\tocsub{I.\quad Works on the Lahu language and people}
\tocsub{II.\quad Works written in Lahu}
\tocsub{III.\quad Works consulted in preparing this \textit{Dictionary}}
\tocline{Plates}{toc:plates}

\clearpage

\input{generated/latex/frontmatter-post-toc.tex}

\clearpage
\pagenumbering{arabic}
\setcounter{page}{1}
\pagestyle{fancy}

\begin{multicols}{2}
\raggedright
\small
\setlength{\dictindent}{1.3em}
\input{generated/latex/lahu.tex}
\end{multicols}

%% Back matter (appendices, bibliography, plate captions) continues the
%% dictionary body's arabic numbering, standard book convention -- see
%% src/process_front_and_back_matter.py for which files/order/why. Same
%% reasoning as the front matter above: back on \pagestyle{plain}, not
%% the dictionary's own fancy collation footer.
\pagestyle{plain}
\input{generated/latex/backmatter.tex}

\end{document}
